% ------------------------------------------------------------
% §4  Random mappings and Flajolet's lens
% ------------------------------------------------------------
\section{Random mappings: Flajolet's telescope}\label{sec:randommaps}

If the working conjecture is ``\SHA{} behaves like a random function,''
then the right response is the scientist's: \emph{a conjecture that
mentions randomness is a bundle of quantitative predictions}. What does a
random function actually look like? Remarkably, this question has an
exact, beautiful answer, developed by Harris in the probabilistic
era~\citep{Harris1960} and perfected by Flajolet and Odlyzko with
generating-function methods~\citep{FlajoletOdlyzko1990}. The answer turns
the vague conjecture into a table of measurable constants.

\subsection{The functional graph of a random map}\label{sec:fungraph}

Let $\varphi$ be a function from a finite set of size $N$ to itself---for
us, eventually, $\varphi = $ (some restriction of) \SHA{}, with
$N = 2^n$. Draw the \emph{functional graph}: a node for each point, an
arrow $x \to \varphi(x)$. Because every node has out-degree exactly one,
the shape is forced: each connected component consists of exactly one
directed cycle, with trees hanging off the cycle nodes
(\cref{fig:rho}). Iterating $\varphi$ from any seed $x_0$ walks a
$\rho$-shaped path: a \emph{tail} of $\lambda$ new points, then a
\emph{cycle} of length $\mu$ repeated forever.

\begin{figure}[ht]
  \centering
  \begin{tikzpicture}[>={Stealth[length=1.8mm]},
      pt/.style={circle, fill, inner sep=1.4pt},
      lbl/.style={font=\scriptsize}]
    % cycle of 8, label inside
    \foreach \i in {0,...,7} {
      \pgfmathsetmacro{\ang}{\i*45}
      \node[pt] (c\i) at ($(4.55,0)+(\ang:1.2)$) {};
    }
    \foreach \i [evaluate=\i as \j using {int(mod(\i+1,8))}] in {0,...,7}
      \draw[->] (c\i) to[bend right=18] (c\j);
    \node[lbl, align=center] at (4.55,0) {cycle\\length $\mu$};
    % horizontal tail entering the leftmost cycle node c4
    \node[pt] (t0) at (0.4,0) {};
    \node[pt] (t1) at (1.4,0) {};
    \node[pt] (t2) at (2.4,0) {};
    \node[lbl, left=1mm of t0] {$x_0$};
    \draw[->] (t0) -- (t1);
    \draw[->] (t1) -- (t2);
    \draw[->] (t2) -- (c4);
    % small trees hanging off cycle nodes, kept clear of everything
    \node[pt] (a1) at ($(c2)+(-0.45,0.75)$) {};
    \node[pt] (a2) at ($(c2)+(0.5,0.8)$) {};
    \node[pt] (a3) at ($(a2)+(0.55,0.65)$) {};
    \draw[->] (a1) -- (c2); \draw[->] (a2) -- (c2); \draw[->] (a3) -- (a2);
    \node[pt] (b1) at ($(c7)+(0.65,-0.6)$) {};
    \node[pt] (b2) at ($(b1)+(0.55,-0.6)$) {};
    \draw[->] (b1) -- (c7); \draw[->] (b2) -- (b1);
    % tail brace
    \draw[decorate, decoration={brace, amplitude=4pt, mirror}]
      (0.4,-0.35) -- (3.1,-0.35)
      node[midway, below=5pt, lbl] {tail, length $\lambda$};
  \end{tikzpicture}
  \caption{One component of a functional graph. Iteration from any seed
  traces a $\rho$: a tail of length $\lambda$ into a cycle of length
  $\mu$. Trees hang off the cycle; components with their unique cycles
  tile the whole space.}
  \label{fig:rho}
\end{figure}

Now let $\varphi$ be chosen \emph{uniformly at random} among all $N^N$
self-maps. Flajolet and Odlyzko computed the expected value of essentially
every natural parameter of this picture; a sample, asymptotically as
$N \to \infty$~\citep{FlajoletOdlyzko1990}:

\begin{center}
\begin{tabular}{@{}llll@{}}
  \toprule
  parameter && expectation & at $N = 2^{256}$ \\
  \midrule
  tail length $\lambda$ (random seed)   && $\sqrt{\pi N / 8}$        & $\approx 2^{127.3}$ \\
  cycle length $\mu$ (random seed)      && $\sqrt{\pi N / 8}$        & $\approx 2^{127.3}$ \\
  rho length $\lambda + \mu$            && $\sqrt{\pi N / 2}$        & $\approx 2^{128.3}$ \\
  number of cyclic points               && $\sqrt{\pi N / 2}$        & $\approx 2^{128.3}$ \\
  number of components                  && $\tfrac12 \ln N$          & $\approx 89$ \\
  image size $|\varphi(\State)|$        && $(1 - e^{-1})\,N$         & $\approx 0.632\,N$ \\
  points with no preimage (leaves)      && $e^{-1}\, N$              & $\approx 0.368\,N$ \\
  largest component                     && $\approx 0.7582\,N$       & --- \\
  largest tree                          && $\approx 0.48\,N$         & --- \\
  \bottomrule
\end{tabular}
\end{center}

\noindent
Pause on the flavor of these results. A ``typical'' random map is nothing
like a random permutation: it is lossy (a third of the space is never hit
at all), it has one giant component swallowing three quarters of
everything, its cycles are tiny compared to $N$ (order $\sqrt{N}$), and
the number of components grows only logarithmically. This is a sharp,
peculiar, recognizable statistical silhouette---a fingerprint of
structurelessness itself.

\subsection{Where the constants come from}\label{sec:singularity}

The route to that table is pure Flajolet, and worth a paragraph even in an
overview, because it connects our discrete object to complex analysis. A
tree rooted at a node whose arrow enters a cycle is a labelled rooted
tree; its exponential generating function $T(z)$ satisfies the functional
equation
\[
  T(z) = z\, e^{T(z)}
\]
(a root, plus a set of subtrees). Components are cycles of trees, and
mappings are sets of components, which by the symbolic calculus gives the
mapping EGF
\[
  M(z) \;=\; \frac{1}{1 - T(z)} .
\]
All the parameters in the table are obtained by marking features in these
equations and then reading off coefficient asymptotics. The engine for the
last step: $T$ has a square-root singularity at $z = e^{-1}$, where
$T(e^{-1}) = 1$---so $M$ blows up there---and \emph{singularity analysis}
transfers the local behavior of a generating function at its dominant
singularity directly into asymptotics of its coefficients. The
half-integer exponents in the singular expansions are what produce all
those $\sqrt{N}$'s; the method is developed at book length
in~\citet{FlajoletSedgewick2009}, with random mappings as a recurring
worked example. For a mathematician, this is the pleasant surprise of the
whole subject: questions about iterating a hash function route through
the local behavior of an analytic function at a branch point.

\subsection{The telescope, pointed at \SHA{}}\label{sec:telescope}

Here is the epistemological payoff. The random-function conjecture of
\cref{sec:history} predicts that if we build from \SHA{} a self-map of an
$N$-point space---most simply
\[
  \varphi_n : \bitstrings^n \to \bitstrings^n, \qquad
  \varphi_n(x) = \text{first } n \text{ bits of } \SHA(x),
\]
for moderate $n$---then the functional graph of $\varphi_n$, an object
\SHA{}'s designers never contemplated and certainly never tuned, should
match every line of the Flajolet--Odlyzko table, not merely in expectation
but in distribution (the limiting distributions are also known). Each
parameter is a falsifiable prediction; jointly they are a
\emph{coordinate system} in which the position of \SHA{} among
pseudorandom objects can be measured rather than argued about. For
contrast, structured maps land visibly elsewhere in these coordinates:
$x \mapsto x \shaplus c$ decomposes into cycles only (no trees: it is a
bijection), $x \mapsto x^2 \bmod p$ has its graph dictated by quadratic
reciprocity, and a one-round toy \SHA{} should sit measurably off the
random silhouette. Where exactly it rejoins the silhouette as rounds are
added is, as far as we know, unmeasured---and eminently measurable
(\cref{sec:projects}).

At full scale ($n = 256$, $N = 2^{256}$) nobody will ever traverse a rho:
the expected rho length is about $2^{128}$ steps. The table is then the
\emph{only} available description of the global geometry of the object the
whole world is using---a telescope pointed at a structure we can predict
in detail but never visit. And the predictions are not idle: they are the
engineering basis of real algorithms. Pollard's rho
method~\citep{Pollard1978} finds collisions and discrete logarithms
precisely by betting that its iteration behaves like a random map (cycle
detection in $O(\sqrt{N})$ steps with $O(1)$ memory); van Oorschot and
Wiener's parallel collision search~\citep{vanOorschotWiener1999} turned
that bet into an industrial discipline (distinguished points, linear
parallel speedup), and the SHAttered attack's budget of $2^{63}$
(\cref{sec:falling}) was costed with exactly this calculus. Every time
such an algorithm's observed running time matches its predicted constant,
a Flajolet statistic of a real hash function has been confirmed to one
more decimal place.

So far the telescope has resolved a \emph{single} application of the
function. Because \SHA{} is so often composed with itself---in password
stretching, hash chains, and Bitcoin's double hash---the dynamics of
\emph{iterating} $\varphi$ form their own coordinate system, and
\cref{sec:iterated} takes it up next: how the image shrinks, where
orbits settle, and the precise sense in which iteration almost, but not
quite, induces a random permutation.

\begin{project}{A functional-graph atlas of truncated \SHA{}}
  {programming; a first probability course}{6--8 weeks (flagship)}
  For $n = 16, 18, \dots, 30$: exhaustively build the functional graph of
  $\varphi_n$, measure every parameter in the table above (and the
  distributions, not just means), and compare with the
  Flajolet--Odlyzko asymptotics \emph{and} with genuinely random maps of
  the same size as a control. Deliverable: an ``atlas'' of plots---theory
  curve, random-map control cloud, \SHA{} data points. Then repeat with
  reduced-round \SHA{} ($1, 2, 4, \dots, 64$ rounds) and watch the data
  walk onto the random silhouette as rounds accumulate. To our knowledge
  no published systematic atlas of this kind exists; undergraduates can
  own it.
\end{project}

\begin{project}{Collision hunting, the van Oorschot--Wiener way}
  {programming; probability}{4--6 weeks}
  Implement Pollard rho with Brent or Floyd cycle detection on
  $\varphi_n$ for $n$ up to ${\sim}48$, then the van Oorschot--Wiener
  distinguished-points parallel method on a lab's worth of machines. Measure the distribution of time-to-collision against the
  theory of \cref{sec:fungraph}. Students reproduce, at toy scale, the
  exact methodology of the SHAttered attack---and learn why $256$ bits is
  chosen so that $\sqrt{N}$ is out of humanity's reach.
\end{project}

\subsection{Coda: counting to three, and the dials of \SHA{}}
\label{sec:josephus3}

Before putting the telescope away, we settle a debt incurred in the
opening paragraph of this paper (\cref{sec:josephus}). The Josephus
circle there counted to \emph{two}, and the survivor map collapsed to a
one-bit rotation. What happens when the circle counts to three? The
answer is known, and it comes---pleasingly---from the same Odlyzko whose
lens this section has been borrowing~\citep{OdlyzkoWilf1991}. The
survivor $J_3(n)$ is delivered by a crisp algorithm~\citep[\S 3.3]{GKP1994}:
iterate $D_0 = 1$, $D_k = \bigl\lceil \tfrac{3}{2} D_{k-1} \bigr\rceil$,
stop at the first $D_k > 2n$, and take $J_3(n) = 3n + 1 - D_k$. Odlyzko
and Wilf proved that these iterates admit an exact closed form,
\[
  D_k \;=\; \Bigl\lfloor K \bigl(\tfrac{3}{2}\bigr)^{k} \Bigr\rfloor
  \quad (k = 0, 1, \dots), \qquad
  K \;=\; 1.62227\,05028\,84767\ldots,
\]
so $J_3$, like $J_2$, ``has a formula.'' But inspect the formula. The
constant $K$ is assembled from the error terms of the very recurrence it
is supposed to shortcut; no independent characterization of it is known,
so the only way anyone can evaluate the closed form is to run the
iteration it was meant to abolish---as the authors say plainly, the
explicit formula ``is not an improvement over the algorithm.'' Nor is
this an artifact of technique. For the general iteration
$f_{n+1} = \lceil \alpha f_n \rceil$ the analogous constant $c(\alpha)$
is a jagged function of $\alpha$, with jump discontinuities at every
integer and at every Josephus point $\alpha = q/(q-1)$; and the paper
closes with a conjecture on the fine distribution of the error terms
that its authors liken to proving that a given real number is normal in
a given base---``likely to be very difficult.'' Inside this family the
solved case is a degenerate point: at $q = 2$ the error terms vanish
identically, $K(2) = 1$, and the machinery collapses back to the bit
rotation of \cref{eq:josephus}.

Dwell on what just happened. We did not change the problem; we turned
one dial by one click, from every-second to every-third, and
analyzability fell off a cliff---from a closed form a child can verify
to a constant nobody can compute independently and a conjecture
calibrated against normality. The collapse that opened this paper was
not a property of decimation in general. It was an accident of the
parameter $q = 2$. In the family, wildness is generic; solvability is
the exception.

\SHA{} deserves to be examined in exactly this light, because its
definition is also a single point in an evident parameter family, with
many more dials: the rotation offsets $(2,13,22)$ and $(6,11,25)$ of
$\Sigma_0$ and $\Sigma_1$, the rotation-and-shift offsets $(7,18;\,3)$
and $(17,19;\,10)$ of $\sigma_0$ and $\sigma_1$, the word length $32$,
the round count $64$, the schedule taps $(2,7,15,16)$, the sixty-four
additive constants $K_t$. Each could have been set otherwise, and asking
what happens as one is perturbed---the analogue of moving $q$---is a
legitimate and largely unexplored experimental axis. In one direction
the family provably degenerates: set the three rotation offsets of
$\Sigma_0$ equal and the map collapses to a single rotation, as
transparent as the Josephus shift; drop to an even number of rotation
summands and invertibility itself dies (\cref{thm:rivest}); move the
word length off a power of two and the parity law is repealed. The
published constants are believed to sit in the opposite, generic regime,
where nothing collapses---but between the degenerate points and the
conjecturally generic bulk, nobody has mapped the terrain. The
functional-graph atlas proposed above thus gains a second axis: measure
the Flajolet coordinates of reduced-round variants not only as rounds
accumulate but as the rotation and shift amounts are turned away from
the standard's values, and ask whether the chosen point is
distinguishable, in any coordinate of the table, from a randomly chosen
neighbor.
